\documentclass[11pt]{article}
\usepackage[a4paper,margin=26mm]{geometry}
\usepackage{amsmath,amssymb,amsthm,hyperref}
\newtheorem{theorem}{Theorem}
\newtheorem{corollary}[theorem]{Corollary}
\newcommand{\Aut}{\operatorname{Aut}}
\title{An orbit bound for primitive geometries}
\author{A correction to the Rudvalis entry of Leemans--Rodrigues}
\date{}
\begin{document}\maketitle
Leemans and Rodrigues classify rank-two primitive coset geometries using
unordered pairs of distinct maximal subgroups, with equivalence under
the automorphism group of the ambient group. Table~1 of their paper
\cite{LR}, on physical page~5, gives $1511$ geometries for the Rudvalis
group. Under precisely this convention, a single conjugacy class of
maximal subgroups already gives at least $50669$. The argument is an
orbit bound; it does not require enumerating these geometries.

\begin{theorem}\label{bound}
Let $G$ be a finite nonabelian simple group with trivial outer
automorphism group. Suppose that $M$ is a maximal subgroup of order $m$,
and put $n=|G|/m$. There are at least
\[
                   \left\lceil\frac{n-1}{2m}\right\rceil
\]
equivalence classes of firm, connected, primitive rank-two coset
geometries whose two parabolic subgroups are conjugate to $M$.
Equivalence here allows interchange of the two types and is induced
by an element of $\Aut(G)$.
\end{theorem}
\begin{proof}
Maximality gives $N_G(M)=M$ or $G$. The latter alternative makes $M$
normal. It is impossible because $G$ is simple and $M$ is a nontrivial
proper subgroup. Thus $M$ has exactly $n$ conjugates. Let $\Omega$ be
this set, and let $X$ be the set of its two-element subsets. Then
$|X|=n(n-1)/2$.

Each $\{H,K\}\in X$ defines a rank-two coset geometry. The right cosets
of $H$ and $K$ are incident when their intersection is nonempty. Every
incident pair is carried to $\{H,K\}$ by right multiplication, so the
action is flag-transitive. Its local valencies are
$[H:H\cap K]$ and $[K:H\cap K]$. Both exceed one: inclusion between
two distinct maximal subgroups is impossible. Hence the geometry is
firm. Moreover $\langle H,K\rangle=G$ by maximality. Paths in the
incidence graph correspond to alternating products from $H$ and $K$,
so this equality proves connectedness. Maximality of both parabolic
subgroups gives primitivity on each type. Simplicity also makes both
coset actions faithful, since their kernels are proper normal subgroups.

The equivalence stipulated in the statement is exactly the orbit
relation of $\Aut(G)$ on $X$. Since $G$ is centreless and has trivial
outer automorphism group, $\Aut(G)=\operatorname{Inn}(G)\cong G$.
Every orbit therefore has at most $|G|$ elements. Dividing $|X|$ by
$|G|=nm$, and rounding upwards, proves the bound. The factor two has
already allowed interchange of $H$ and $K$; it must not be charged again.
\end{proof}

\begin{corollary}\label{ru}
There are at least $50669$ primitive rank-two coset geometries for the
Rudvalis group, with equivalence as in \cite{LR}. In particular, the
entry $1511$ in its Table~1 is false under its stated convention.
\end{corollary}
\begin{proof}
We retain the following classical group facts from the Atlas
\cite{Atlas}: the Rudvalis group is nonabelian simple, has order
$145926144000$, has trivial outer automorphism group, and contains a
maximal subgroup $(5{:}4)\times A_5$. These are inputs, not new results
of this note. The subgroup has order $20\cdot60=1200$, and hence
\[
 n=\frac{145926144000}{1200}=121605120.
\]
Theorem~\ref{bound} now gives
\[
 \left\lceil\frac{121605119}{2400}\right\rceil=50669>1511.
\]
Both defining subgroups belong to just this one conjugacy class, so
including all maximal subgroups can only increase the number.
\end{proof}

There is also a useful exact formula underlying the bound. It separates
the counting problem from the construction of representatives.
\begin{theorem}\label{character}
Let a finite group $A$ act on a finite set $\Omega$, and let
$\psi(a)=|\operatorname{Fix}_{\Omega}(a)|$. The number of $A$-orbits
on two-element subsets is
\[
 \frac{1}{2|A|}\sum_{a\in A}
             \bigl(\psi(a)^2-2\psi(a)+\psi(a^2)\bigr).       \tag{1}
\]
For a nonabelian simple group $G$, take $A=\Aut(G)$ and $\Omega$ to
be all maximal subgroups of $G$. Formula~(1) then counts exactly the
primitive rank-two geometries with the equivalence of \cite{LR}.
\end{theorem}
\begin{proof}
A two-element subset fixed by $a$ either consists of two individually
fixed points, or forms a two-cycle of $a$. There are
$\binom{\psi(a)}2$ subsets of the first kind. Of the $\psi(a^2)$ points
fixed by $a^2$, exactly $\psi(a)$ are already fixed by $a$; the remainder
occur in two-cycles. Thus the number of fixed subsets is
$(\psi(a)^2-2\psi(a)+\psi(a^2))/2$. Burnside's orbit-counting lemma
gives (1). The last assertion follows from the same coset and maximality
argument as in Theorem~\ref{bound}.
\end{proof}

The bound in Corollary~\ref{ru} is the certified numerical conclusion.
It is not a corrected exact total or a construction of every geometry.
The other eleven rows of Table~1, the printed orbit-length tables,
the designs and the external census data require separate verification.
No claim about them follows merely from this contradiction. Nor does
this note establish that the fixed edition has no later correction.

The source supplies the geometry problem and equivalence convention;
the Atlas supplies the stated group data. Orbit--stabilizer and
Burnside's lemma are classical. Our contribution here is their
application to the source's numerical assertion and the explicit
separation between counting orbits and constructing representatives.
The accompanying integer certificate checks the displayed arithmetic;
the general argument above supplies the proof. We claim neither a
newly solved open conjecture nor a compression factor in this note.

\begin{thebibliography}{2}
\bibitem{LR} D.~Leemans and B.~G.~Rodrigues,
\emph{On primitive geometries of rank two},
The Art of Discrete and Applied Mathematics \textbf{2} (2019), P1.08.
\href{https://doi.org/10.26493/2590-9770.1301.165}{doi:10.26493/2590-9770.1301.165}.
\bibitem{Atlas} R.~A.~Wilson, R.~A.~Parker and J.~N.~Bray,
\emph{Atlas of Finite Group Representations: Rudvalis group Ru},
group data and maximal subgroup list, accessed 14 October 2026.
\url{https://brauer.maths.qmul.ac.uk/Atlas/spor/Ru/}.
\end{thebibliography}
\end{document}
